\section{Large Sets} \label{sec:filters}
\subsection{A new point of view}
In order to build the hyperreal numbers we need one new key element: a new equivalence relation that lets us compare sequences of real numbers. We could make use of the equivalence relation used in the construction of $\Rs$ from $\mathbb{Q}$, which is the eventually $\varepsilon$-close equivalence relation used by \cite{tao2022analysis}. However, despite being a sensible approach and despite it having been a bullet-proof criteria upon which to construct the real numbers, it wouldn't do us any good. The reason is that, when deploying the hyperreal numbers and making use of non-standard analysis, one of the true powers lies in the fact that we can regard the sequence $(\frac{1}{n})_{n=0}^\infty$ as an element different from $0$. We could make use of such sequences to define elements that, while always being strictly different from zero, keep approaching $0$ and thus become arbitrarily small, or \textit{infinitesimal}. This is clearly not possible within the framework of standard real analysis, where these two elements are regarded to be the same since for every $\varepsilon>0$ there exists $N>0$ such that $\abs{a_n - a_m} \leq \varepsilon$ for all $n, m > N$. Nonetheless, there is no term in this sequence, no matter how far in the sequence you go, that will actually be equal to zero.
According to our opinion, this is one of the greatest limitations of this equivalence criterion of \textit{closeness}. Thus, we try to forget this relation of closeness and we aim for a more permissive equivalence relation.
This new relation has something to do with the question: when is a set large?

\subsection{Largeness}
Let $r=(r_1,r_2,r_3,...)$ and $s=(s_1,s_2,s_3,...)$ be real-valued sequences. We will say that $r$ and $s$ are equivalent if they agree at a "large" number of places, i.e., if their \textit{agreement set}
\begin{equation*}
    E_{rs} = \{n:r_n=s_n\}
\end{equation*}
is \textit{large}. This definition is obviously vague for the simple reason that we have not yet defined what large means. Bear with us. Regardless of the meaning of "large", we would like some properties to be satisfied:
\begin{itemize}
    \item $\Ns$ must be large.
    \item Equivalence must be a transitive relation. If $A$ and $B$ are large sets, and $A \cap B \subseteq C$ then $C$ is large.
    \item The empty set $\emptyset$ is not large. 
    \item For any subset $A$ of $\Ns$, one of $A$ and $\Ns \setminus A$ is large.
\end{itemize}

Is there such a notion of largeness?

\subsection{Filters}
Let $I$ be a nonempty set. The \textit{power set} of $I$ is the set
\begin{equation*}
    \mathcal{P} = \{A : A \subseteq I\},
\end{equation*}
the set of all subsets of $I$. 

\defn{A \textit{filter} on $I$ is a nonempty collection $\mathcal{F} \subseteq \cal{P}$$(I)$ of subsets of $I$ satisfying the following axioms:
\begin{itemize}
    \item Intersections: if $A, B \in \cal{F}$, then $A \cap B \in \cal{F}$.
    \item Supersets: if $A \in \cal{F}$ and $A \subseteq B \subseteq I$, then $B \in \cal{F}$.
\end{itemize}}

\prop{A filter $\cal{F}$ contains the empty set $\emptyset$ iff $\mathcal{F}=\mathcal{P}(I)$, i.e,
\begin{equation*}
    \emptyset \in \call{F} \iff \call{F}=\call{P}(I)
\end{equation*}
\begin{proof}
    ($\Leftarrow$) This direction follows trivially: we know that $\emptyset \in \call{P}(I)$ and $\call{F}=\call{P}(I)$, so it follows that $\emptyset \in \call{F}$.

    \noin Let us now prove the opposite direction ($\Rightarrow$). Let $\emptyset \in \call{F}$. We want to prove $\call{F}=\call{P}(I)$. To do so we must prove that $\call{F} \subseteq \call{P}(I)$ and $\call{P}(I)\subseteq \call{F}$. Let's start with the easiest part: 
    \begin{enumerate}
        \item by definition, a filter $\call{F}$ on $I$ is a nonempty collection $\call{F} \subseteq \call{P}(I)$ of subsets of $I$, so  $\call{F} \subseteq \call{P}(I)$ follows trivially.
        \item Let us prove that $\call{P}(I)\subseteq \call{F}$. Let $B \in \call{P}(I)$, so we know that $B \subseteq I$. We know that $\emptyset \in \call{F}$ and that $\emptyset \subseteq B \subseteq I$. By the superset axiom, we must have that $B \in \call{F}$. Thus $\call{P}(I)\subseteq \call{F}$.
    \end{enumerate} 
    We proved both sides of the $\subseteq$ relation, thus $\call{P}(I) = \call{F}$. We also proved both implications thus the proposition holds. 
\end{proof}

\defn[Proper and Ultrafilter]{$\cal{F}$ is \textit{proper} if $\emptyset \notin \cal{F}$. An \textit{ultrafilter} is a proper filter that satisfies:
\begin{equation*}
\text{for any } A \subseteq I,\ \text{either } A \in \mathcal{F}
\text{ or } A^c \in \mathcal{F},
\qquad \text{where } A^c = I \setminus A.
\end{equation*}
}

\subsection{Examples of Filters}
\subsubsection{Principal ultrafilters} \label{eg:principal ultrafilter generated by i}
\begin{defn}[Principal ultrafilter generated by \(i\)]

Let \(i \in I\). Then
\[
\mathcal{F}^{i}
=
\left\{
A \subseteq I : i \in A
\right\}
\]
is an ultrafilter, called the \emph{principal ultrafilter generated by}
\(i\).

If \(I\) is finite, then every ultrafilter on \(I\) is of the form
\(\mathcal{F}^{i}\) for some \(i \in I\), and so is principal.

Think of an ultrafilter as a rule that declares certain subsets of \(I\) to
be ``large.'' A principal ultrafilter generated by \(i\) uses an extremely
simple rule:
\begin{equation*}
\boxed{\text{A subset is large exactly when it contains } i.}
\end{equation*}

Formally,
\[
\mathcal{F}^i = \{A \subseteq I : i \in A\}.
\]
Here, \(i\) is one fixed, distinguished element of \(I\).
\end{defn}

\begin{example}[A concrete example]
Let \(I = \{1,2,3,4\}\) and choose \(i=2\). Then
\[
\mathcal{F}^2 = \{A \subseteq I : 2 \in A\}.
\]
So \(\mathcal{F}^2\) consists of all subsets containing \(2\):
\begin{equation*}
\mathcal{F}^2 =
\bigl\{
\{2\},
\{1,2\},
\{2,3\},
\{2,4\},
\{1,2,3\},
\{1,2,4\},
\{2,3,4\},
\{1,2,3,4\}
\bigr\}.
\end{equation*}

It does not contain sets such as \(\emptyset\), \(\{1\}\), and \(\{3,4\}\),
because none of those contains \(2\).

Visually, the power set is divided into two equal groups:
\[
\begin{array}{c|c}
\text{Contains } 2 & \text{Does not contain } 2 \\
\hline
\text{belongs to } \mathcal{F}^2 & \text{does not belong to } \mathcal{F}^2
\end{array}
\]
\end{example}

\begin{prop}
The collection \(\mathcal{F}^i\) is a proper filter on \(I\).
\end{prop}

\begin{proof}
We check the two filter properties.

Suppose \(A,B \in \mathcal{F}^i\). By definition, \(i \in A\) and
\(i \in B\). Therefore, \(i \in A \cap B\). Hence,
\[
A \cap B \in \mathcal{F}^i.
\]

Intuitively, if two sets both contain the distinguished point \(i\), then
their intersection still contains \(i\).

Now suppose \(A \in \mathcal{F}^i\) and \(A \subseteq B \subseteq I\).
Since \(i \in A\) and \(A \subseteq B\), we must have \(i \in B\).
Therefore,
\[
B \in \mathcal{F}^i.
\]

Intuitively, once a set contains \(i\), making the set larger cannot remove
\(i\).

Also,
\[
\emptyset \notin \mathcal{F}^i,
\]
because \(\emptyset\) contains no elements. Thus \(\mathcal{F}^i\) is a
proper filter.
\end{proof}

\begin{prop}
The collection \(\mathcal{F}^i\) is an ultrafilter.
\end{prop}

\begin{proof}
For every subset \(A \subseteq I\), there are exactly two possibilities:
\(i \in A\) or \(i \notin A\).

If \(i \in A\), then \(A \in \mathcal{F}^i\). If \(i \notin A\), then
\(i\) must lie in the complement \(A^c = I \setminus A\). Therefore,
\[
A^c \in \mathcal{F}^i.
\]

So for every \(A \subseteq I\),
\[
A \in \mathcal{F}^i
\quad \text{or} \quad
A^c \in \mathcal{F}^i.
\]

In fact, exactly one belongs to the ultrafilter: \(A\) and \(A^c\) cannot
both contain \(i\). That is precisely the ultrafilter property.
\end{proof}

\paragraph{The ultrafilter \(\mathcal{F}^i\) is principal.}
The singleton \(\{i\}\) is the smallest set in \(\mathcal{F}^i\). Every
other member of the filter is a superset of \(\{i\}\):
\[
A \in \mathcal{F}^i
\quad \Longleftrightarrow \quad
\{i\} \subseteq A.
\]

Thus the whole ultrafilter is generated by the one-element set \(\{i\}\):
\[
\mathcal{F}^i = \{A \subseteq I : \{i\} \subseteq A\}.
\]

The point \(i\) acts as the ``principal'' or central point determining every
membership decision.

\subsubsection{The cofinite filter}

\begin{defn}[Cofinite filter]
The collection
\[
\mathcal{F}^{\mathrm{co}}
=
\left\{
A \subseteq I : I \setminus A \text{ is finite}
\right\}
\]
is the \emph{cofinite} filter on \(I\), and is proper iff \(I\) is infinite.
\(\mathcal{F}^{\mathrm{co}}\) is not an ultrafilter.
\end{defn}

\begin{prop}
The collection \(\mathcal{F}^{\mathrm{co}}\) is a filter.
\end{prop}

\begin{proof}
Let \(A, B \in \call{F}^{\mathrm{co}}\). We need to show that
\(I \setminus (A \cap B)\) is finite. We see that
\begin{align*}
    I \setminus (A \cap B)
    & = I \cap (A \cap B)^c \\
    & = I \cap (A^c \cup B^c) \\
    & = (I \cap A^c) \cup (I \cap B^c) \\
    & = (I \setminus A) \cup (I \setminus B).
\end{align*}
Since the union of finite sets is itself finite, we have that
\(I \setminus (A \cap B) \in \call{F}^{\mathrm{co}}\).

Now let \(A \in \call{F}^{\mathrm{co}}\) and \(A \subseteq B \subseteq I\).
We can prove very easily the following relation,
\[
A \subseteq B \Rightarrow B^c \subseteq A^c.
\]
Since \(A^c\) is finite, every subset of a finite set is finite. Thus
\(B^c\) is also finite, and so
\[
B^c=I\setminus B \in \call{F}^{\mathrm{co}}.
\]
\end{proof}

\begin{prop}
The cofinite filter \(\call{F}^{\mathrm{co}}\) is proper if and only if
\(I\) is infinite.
\end{prop}

\begin{proof}
\noindent\((\Rightarrow)\)
Assume \(\call{F}^{\mathrm{co}}\) is proper, so
\(\emptyset \notin \call{F}^{\mathrm{co}}\). For the sake of contradiction,
assume that \(I\) is finite. Since \(I\) is
finite,
\[
\emptyset = I^c \in \call{F}^{\mathrm{co}},
\]
which is a contradiction. Thus \(I\) must be infinite.

\noindent\((\Leftarrow)\)
Assume \(I\) is infinite. For the sake of contradiction, assume
\(\call{F}^{\mathrm{co}}\) is not proper, thus
\(\emptyset \in \call{F}^{\mathrm{co}}\). Since
\(\emptyset \in \call{F}^{\mathrm{co}}\),
\[
I \setminus \emptyset=I
\]
must be finite. But then this contradicts the assumption. Thus
\(\call{F}^{\mathrm{co}}\) must be proper.
\end{proof}

\begin{prop}
Assume \(I\) is infinite. Then \(\call{F}^{\mathrm{co}}\) is not an
ultrafilter.
\end{prop}

\begin{proof}
Choose pairwise distinct elements \(i_1, i_2,... \in I\), and set
\[
A=\{i_2, i_4,i_6,...\}.
\]
Then \(A\) is infinite. Its complement \(I \setminus A\) is also infinite
because it contains the infinite set
\[
\{i_1, i_3, i_5,...\}.
\]

Since \(I \setminus A\) is infinite we must have that
\[
A \notin \call{F}^{\mathrm{co}},
\]
by definition. But we also know that
\[
I \setminus (I \setminus A) = A
\]
is also infinite, which implies that
\[
I \setminus A \notin \call{F}^{\mathrm{co}}.
\]

Thus \(\call{F}^{\mathrm{co}}\) does not satisfy the ultrafilter property,
and so it is not an ultrafilter.
\end{proof}

\subsubsection{Filters generated by collections}

\begin{defn}[Filter generated by \(\call{H}\)]
If \(\emptyset \neq \mathcal{H} \subseteq \call{P}(I)\), then the
\emph{filter generated by} \(\call{H}\), i.e., the smallest filter on \(I\)
including \(\call{H}\), is the collection
\[
\call{F}^\call{H}
=
\left\{
A \subseteq I :
B_1 \cap \cdots B_n \subseteq A
\text{ for some } n \text{ and some } B_i \in \call{H}
\right\}.
\]

If \(\call{H}\) has a single member \(B\), then
\[
\call{F}^\call{H}
=
\{ A \subseteq I : B \subseteq A \},
\]
which is called the \emph{principal filter generated by} \(B\). The
ultrafilter \(\call{F}^i\) of
Example \ref{eg:principal ultrafilter generated by i} is the special case of
this when \(B=\{i\}\).

This is the standard construction that turns an arbitrary collection of
subsets \(\call{H}\) into the smallest filter containing them. It does so by
enforcing exactly the two filter rules: finite intersections, and include
every superset of those intersections. Intuitively, think of each
\(B \in \call{H}\) as a requirement that ``\(B\) must count as large''. A
filter must then accept intersections of such large elements, and if a set
contains a large set, then it too must be large.
\end{defn}

\begin{prop}[\textit{this used to be an exercise}]
The collection \(\mathcal{F}^{\mathcal H}\) contains \(\mathcal H\) and is a
filter.
\end{prop}

\begin{proof}
To show it, we must check that it
contains \(\mathcal H\), is closed under supersets, and is closed under finite
intersections. Recall that
\[
\mathcal{F}^{\mathcal H}
=
\left\{
A \subseteq I :
B_1 \cap \cdots \cap B_n \subseteq A
\text{ for some } n \geq 1
\text{ and } B_1,\ldots,B_n \in \mathcal H
\right\}.
\]

Let \(B \in \mathcal H\). We may take just one set from \(\mathcal H\): set
\(n=1\) and \(B_1=B\). Then \(B_1 \subseteq B\), so by the definition of the
generated filter,
\[
B \in \mathcal{F}^{\mathcal H}.
\]

(Superset axiom). Suppose \(A \in \mathcal{F}^{\mathcal H}\). Then, by definition, there are
\(B_1,\ldots,B_n \in \mathcal H\) such that
\[
B_1 \cap \cdots \cap B_n \subseteq A.
\]

Now suppose that \(A \subseteq C \subseteq I\). By transitivity of
inclusion,
\[
B_1 \cap \cdots \cap B_n \subseteq A \subseteq C.
\]

Thus \(C\) contains the same finite intersection of members of
\(\mathcal H\). Therefore,
\[
C \in \mathcal{F}^{\mathcal H}.
\]

(Intersection axiom). Let \(A,C \in \mathcal{F}^{\mathcal H}\). Then there are finite collections
\(B_1,\ldots,B_m \in \mathcal H\) and
\(D_1,\ldots,D_k \in \mathcal H\) such that
\[
B_1 \cap \cdots \cap B_m \subseteq A,
\qquad
D_1 \cap \cdots \cap D_k \subseteq C.
\]

Any point lying in every \(B_i\) and every \(D_j\) lies in both \(A\) and
\(C\). Hence
\[
B_1 \cap \cdots \cap B_m
\cap
D_1 \cap \cdots \cap D_k
\subseteq
A \cap C.
\]

The set on the left is still a \emph{finite intersection} of members of
\(\mathcal H\). Therefore,
\[
A \cap C \in \mathcal{F}^{\mathcal H}.
\]

Intuitively, \(A\) meets one finite collection of requirements from
\(\mathcal H\), while \(C\) meets another. Their intersection \(A \cap C\)
meets both collections of requirements simultaneously.
\end{proof}

\begin{prop}[\textit{this used to be an exercise}]
The collection \(\mathcal{F}^{\mathcal H}\) is the smallest filter including
\(\call{H}\).
\end{prop}

\begin{proof}
Let \(\call{G}\) be any filter containing \(\call{H}\), so we have
\[
\call{H} \subseteq \call{G}.
\]
Let \(B_1,...,B_n \in \call{H}\), then
\[
B_1,...,B_n \in \call{G}.
\]
Thus, by the intersection axiom,
\[
B_1 \cap \cdots \cap B_n \in \call{G}.
\]

Now assume that
\[
B_1 \cap \cdots \cap B_n \subseteq A.
\]
By the superset axiom,
\[
A \in \call{G}.
\]

By definition, we also have that
\[
A \in \mathcal{F}^{\mathcal H}.
\]

So if \(\call{G}\) is any filter containing \(\call{H}\), the members of
\(\mathcal{F}^{\mathcal H}\) are forced to be members of \(\call{G}\). But
let us now prove that \(\mathcal{F}^{\mathcal H}\) is actually the smallest
filter containing \(\call{H}\).

For the sake of contradiction assume that \(\call{G}\) is actually the
smallest:
\[
\call{G} \subsetneq \mathcal{F}^{\mathcal H}.
\]
Thus, there exists a set
\[
C \in \mathcal{F}^{\mathcal H}
\]
such that
\[
C \notin \call{G}.
\]

Since \(C \in \mathcal{F}^{\mathcal H}\), we must have that
\[
B_1 \cap \cdots \cap B_n \subseteq C,
\]
for \(B_1,...,B_n \in \call{H}\). But we know that
\[
B_1,...,B_n \in \call{G}
\]
and thus by the superset axiom, we must have that
\[
C \in \call{G},
\]
which is a contradiction. Thus \(\mathcal{F}^{\mathcal H}\) is the smallest
filter including \(\call{H}\).
\end{proof}

\prop{Let \(I\) be a nonempty set and let \(A \subseteq I\). Define
\[
\call{F}^{\{A\}}
=
\{B \subseteq I : A \subseteq B\}.
\]
Then \(\call{F}^{\{A\}}\) is proper if and only if \(A \neq \varnothing\).}

\begin{proof}
Recall that \(\call{F}^{\{A\}}\) is proper if and only if
\(\varnothing \notin \call{F}^{\{A\}}\).

\call{F}irst, assume that \(\call{F}^{\{A\}}\) is proper. Suppose, for
contradiction, that \(A=\varnothing\). Then
\[
A=\varnothing\subseteq\varnothing.
\]
Hence, by the definition of \(\call{F}^{\{A\}}\), we have
\[
\varnothing\in \call{F}^{\{A\}},
\]
which contradicts the assumption that \(\call{F}^{\{A\}}\) is proper.
Therefore, \(A\neq\varnothing\).

Conversely, assume that \(A\neq\varnothing\). Suppose, for
contradiction, that
\[
\varnothing\in \call{F}^{\{A\}}.
\]
By the definition of \(\call{F}^{\{A\}}\), this implies that
\[
A\subseteq\varnothing.
\]
Since the empty set has no elements, the inclusion
\(A\subseteq\varnothing\) implies that \(A=\varnothing\). This
contradicts the assumption that \(A\neq\varnothing\). Therefore,
\[
\varnothing\notin \call{F}^{\{A\}},
\]
and hence \(\call{F}^{\{A\}}\) is proper.
\end{proof}

\subsubsection{Unions of chains of filters}
\begin{prop}[Union of a linearly ordered collection of filters] \label{prop:union of a linearly ordered collection of filters is a filter}
Let \(I\) be a nonempty set, let \(X\neq\emptyset\), and let
\[
\{\mathcal{F}_x : x\in X\}
\]
be a collection of filters on \(I\) that is linearly ordered by set
inclusion. That is, for every \(x,y\in X\), either
\[
\mathcal{F}_x\subseteq\mathcal{F}_y
\qquad\text{or}\qquad
\mathcal{F}_y\subseteq\mathcal{F}_x.
\]
Then
\[
\bigcup_{x\in X}\mathcal{F}_x
=
\left\{
A\subseteq I :
A\in\mathcal{F}_x
\text{ for some }x\in X
\right\}
\]
is a filter on \(I\).
\end{prop}

\begin{proof}
Define
\[
\mathcal{F}
=
\bigcup_{x\in X}\mathcal{F}_x.
\]
We prove that \(\mathcal{F}\) is a filter on \(I\).

First, every member of \(\mathcal{F}\) is a subset of \(I\). Indeed, if
\(A\in\mathcal{F}\), then by the definition of the union there exists
\(x\in X\) such that
\[
A\in\mathcal{F}_x.
\]
Since \(\mathcal{F}_x\) is a filter on \(I\), we have
\[
\mathcal{F}_x\subseteq\mathcal{P}(I),
\]
and hence \(A\subseteq I\). Therefore,
\[
\mathcal{F}\subseteq\mathcal{P}(I).
\]

We must also show that \(\mathcal{F}\) is nonempty. Since \(X\neq\emptyset\),
we may choose \(x_0\in X\). Because \(\mathcal{F}_{x_0}\) is a filter, it is
nonempty. In particular,
\[
I\in\mathcal{F}_{x_0}.
\]
Since
\[
\mathcal{F}_{x_0}\subseteq
\bigcup_{x\in X}\mathcal{F}_x
=
\mathcal{F},
\]
we obtain
\[
I\in\mathcal{F}.
\]
Thus \(\mathcal{F}\neq\emptyset\).

We now verify the intersection axiom. Let
\[
A,B\in\mathcal{F}.
\]
By the definition of \(\mathcal{F}\), there exist \(x,y\in X\) such that
\[
A\in\mathcal{F}_x
\qquad\text{and}\qquad
B\in\mathcal{F}_y.
\]
Since the collection is linearly ordered by inclusion, either
\[
\mathcal{F}_x\subseteq\mathcal{F}_y
\]
or
\[
\mathcal{F}_y\subseteq\mathcal{F}_x.
\]

Suppose first that
\[
\mathcal{F}_x\subseteq\mathcal{F}_y.
\]
Then \(A\in\mathcal{F}_y\), and we already know that
\(B\in\mathcal{F}_y\). Since \(\mathcal{F}_y\) is a filter,
\[
A\cap B\in\mathcal{F}_y.
\]
Therefore,
\[
A\cap B\in\mathcal{F}.
\]

The other case is analogous. If
\[
\mathcal{F}_y\subseteq\mathcal{F}_x,
\]
then both \(A\) and \(B\) belong to \(\mathcal{F}_x\), so
\[
A\cap B\in\mathcal{F}_x\subseteq\mathcal{F}.
\]
Thus \(\mathcal{F}\) is closed under finite intersections.

Finally, we verify the superset axiom. Let
\[
A\in\mathcal{F}
\qquad\text{and}\qquad
A\subseteq B\subseteq I.
\]
Since \(A\in\mathcal{F}\), there exists \(x\in X\) such that
\[
A\in\mathcal{F}_x.
\]
Because \(\mathcal{F}_x\) is a filter and
\[
A\subseteq B\subseteq I,
\]
the superset axiom for \(\mathcal{F}_x\) gives
\[
B\in\mathcal{F}_x.
\]
Consequently,
\[
B\in\mathcal{F}.
\]

Therefore, \(\mathcal{F}\) is a nonempty subcollection of
\(\mathcal{P}(I)\) that is closed under finite intersections and supersets.
Hence
\[
\bigcup_{x\in X}\mathcal{F}_x
\]
is a filter on \(I\).
\end{proof}

\subsection{Facts About Filters}
\prop{If $\call{F} \subseteq \call{P}(I)$ satisfies the superset axiom, then
$$
\call{F} \neq \emptyset \iff I \in \call{F}
$$}
\begin{proof}
    ($\Rightarrow$). Assume $\call{F} \neq \emptyset$. Let $A \in \call{F}$, since $\call{F}$ is nonempty. Then $A \subseteq I$. By the superset axiom, $I \in \call{F}$.

    \noin ($\Leftarrow$). Assume $I \in \call{F}$. Then $\call{F}$ is nonempty. So $\call{F} \neq \emptyset$.
\end{proof}

\prop{An ultrafilter \(\mathcal{F}\) satisfies
\begin{enumerate}
    \item \(A \cap B \in \mathcal{F} \iff A \in \mathcal{F}
    \text{ and } B \in \mathcal{F}\).
    
    \item \(A \cup B \in \mathcal{F} \iff A \in \mathcal{F}
    \text{ or } B \in \mathcal{F}\).
    
    \item \(A^c \in \mathcal{F} \iff A \notin \mathcal{F}\).
\end{enumerate}
}
\begin{proof}
    \begin{enumerate}
        \item ($\Rightarrow$). Say $A \cap B \in \mathcal{F}$. $A \cap B \subseteq A, B \subseteq I$, thus $A, B \in \call{F}$. 
        \newline
        \noin ($\Leftarrow$). Follows by definition. 
        \item ($\Rightarrow$). Say $A \cup B \in \mathcal{F}$. Assume for contradiction that $A \notin \call{F}$ and $B \notin \call{F}$. Then $A^c \in \call{F}$ and $B^c \in \call{F}$. By the intersection axiom, $A^c\cap B^c \in \call{F}$. But $ (A\cup B)^c=A^c\cap B^c \in \call{F}$, which is a contradiction, since $\call{F}$ is an ultrafilter. 
        \newline
        \noin ($\Leftarrow$). Say $A \in \call{F}$ or $B \in \call{F}$. (We now assume $A \in \call{F}$ but the proof works in the same way if you assume $B \in \call{F}$). We must have $A \subseteq A \cup B \subseteq I$. By the superset axiom, $A \cup B \in \call{F}$.
        \item ($\Rightarrow$). Let $A^c \in \call{F}$. Suppose, for the sake of contradiction, that $A \in \call{F}$. Then, by the intersection axiom, $\emptyset=A\cap A^c\in \call{F}$, which is a contradiction because $\call{F}$ is a proper filter, and so it doesn't contain the empty set.
        \newline
        \noin ($\Leftarrow$). Let $A \notin \call{F}$. If we assume that also $A^c \notin \call{F}$ we obtain a contradiction because the ultrafilter property of $\call{F}$ guarantees that either $A \in \call{F}$ or $A^c \in \call{F}$. 

        (This third statement is the most intuitive because it follows straightly from the definition of ultrafilter)
    \end{enumerate}
\end{proof}

\prop{Let $\call{F}$ be an ultrafilter and $\{A_1,..., A_n\}$ a finite collection of pairwise disjoint ($A_i \cap A_j =\emptyset$) sets such that
$$
A_1 \cup \cdots \cup A_n \in \call{F}.
$$
Then $A_i \in \call{F}$ for \textit{exactly one} $i$ such that $1 \leq i \leq n$.} \label{prop: in union of sets exactly one belongs to F}
\begin{proof}
Let $\mathcal{F}$ be an ultrafilter on $I$, and let
$A_1,\dots,A_n \subseteq I$ be pairwise disjoint. Thus,
\[
A_i \cap A_j = \emptyset
\qquad \text{whenever } i \neq j.
\]
Assume that
\[
A_1 \cup \cdots \cup A_n \in \mathcal{F}.
\]
We prove that there exists exactly one index
$i \in \{1,\dots,n\}$ such that $A_i \in \mathcal{F}$.

First, we prove existence. Suppose, for contradiction, that
\[
A_i \notin \mathcal{F}
\qquad \text{for every } i=1,\dots,n.
\]
Since $\mathcal{F}$ is an ultrafilter, for each $i$ we have
\[
A_i^c \in \mathcal{F}.
\]
Because $\mathcal{F}$ is closed under finite intersections,
\[
A_1^c \cap \cdots \cap A_n^c \in \mathcal{F}.
\]
By De Morgan's law,
\[
A_1^c \cap \cdots \cap A_n^c
=
\left(A_1 \cup \cdots \cup A_n\right)^c.
\]
Hence,
\[
\left(A_1 \cup \cdots \cup A_n\right)^c
\in \mathcal{F}.
\]
Together with the assumption
\[
A_1 \cup \cdots \cup A_n \in \mathcal{F},
\]
closure under intersections gives
\[
\left(A_1 \cup \cdots \cup A_n\right)
\cap
\left(A_1 \cup \cdots \cup A_n\right)^c
\in \mathcal{F}.
\]
But
\[
\left(A_1 \cup \cdots \cup A_n\right)
\cap
\left(A_1 \cup \cdots \cup A_n\right)^c
=
\emptyset,
\]
so $\emptyset \in \mathcal{F}$, contradicting the fact that
$\mathcal{F}$ is proper. Therefore, there exists at least one
$i \in \{1,\dots,n\}$ such that
\[
A_i \in \mathcal{F}.
\]

Now we prove uniqueness. Suppose that there are distinct indices
$i,j \in \{1,\dots,n\}$ such that
\[
A_i \in \mathcal{F}
\qquad \text{and} \qquad
A_j \in \mathcal{F}.
\]
Since $\mathcal{F}$ is closed under intersections,
\[
A_i \cap A_j \in \mathcal{F}.
\]
However, the sets are pairwise disjoint, so
\[
A_i \cap A_j = \emptyset.
\]
Thus $\emptyset \in \mathcal{F}$, again contradicting the
properness of $\mathcal{F}$.

Therefore, there is exactly one index
$i \in \{1,\dots,n\}$ such that
\[
A_i \in \mathcal{F}.
\]
\end{proof}

\prop{If an ultrafilter contains a finite set, then it contains a
one-element set and is principal. Hence every nonprincipal ultrafilter
contains all cofinite sets.} \label{prop:finite-member-implies-principal}

\begin{proof}
Let $\mathcal{F}$ be an ultrafilter on $I$, and suppose that
$A \in \mathcal{F}$ is finite.

Since $\mathcal{F}$ is proper, we have $A \neq \emptyset$. Therefore,
there exist distinct elements $a_1,\ldots,a_n \in I$ such that
\[
A=\{a_1,\ldots,a_n\}
=
\{a_1\}\cup\cdots\cup\{a_n\}.
\]
The sets $\{a_1\},\ldots,\{a_n\}$ are pairwise disjoint, and their union
belongs to $\mathcal{F}$. By the preceding proposition on pairwise
disjoint finite unions, exactly one of these singleton sets belongs to
$\mathcal{F}$. Thus, for some $i \in \{1,\ldots,n\}$,
\[
\{a_i\}\in\mathcal{F}.
\]

We now show that $\mathcal{F}$ is the principal ultrafilter generated by
$a_i$. We claim that, for every $B\subseteq I$,
\[
B\in\mathcal{F}
\iff
a_i\in B.
\]

First, suppose that $a_i\in B$. Then
\[
\{a_i\}\subseteq B\subseteq I.
\]
Since $\{a_i\}\in\mathcal{F}$, the superset axiom implies that
\[
B\in\mathcal{F}.
\]

Conversely, suppose that $B\in\mathcal{F}$. If $a_i\notin B$, then
$a_i\in B^c$. Hence
\[
\{a_i\}\subseteq B^c\subseteq I.
\]
Since $\{a_i\}\in\mathcal{F}$, the superset axiom gives
\[
B^c\in\mathcal{F}.
\]
But then both $B$ and $B^c$ belong to $\mathcal{F}$, so closure under
intersections gives
\[
\emptyset=B\cap B^c\in\mathcal{F},
\]
contradicting the fact that $\mathcal{F}$ is proper. Therefore,
$a_i\in B$.

We have proved that
\[
\mathcal{F}
=
\{B\subseteq I:a_i\in B\}
=
\mathcal{F}^{a_i}.
\]
Therefore $\mathcal{F}$ is principal.

Finally, suppose that $\mathcal{F}$ is nonprincipal, and let $C\subseteq I$
be cofinite. We prove that $C\in\mathcal{F}$. Assume instead that
\[
C\notin\mathcal{F}.
\]
Since $\mathcal{F}$ is an ultrafilter,
\[
C^c\in\mathcal{F}.
\]
But $C^c$ is finite, since $C$ is cofinite. By the first part of the
proposition, $\mathcal{F}$ must then be principal, contradicting the
assumption that $\mathcal{F}$ is nonprincipal. Hence
\[
C\in\mathcal{F}.
\]

Thus every nonprincipal ultrafilter contains every cofinite subset of $I$.
\end{proof}

\subsubsection{Finite intersection property}

\begin{defn}[Finite intersection property]
A collection \(\mathcal{H} \subseteq \mathcal{P}(I)\) has the
\emph{finite intersection property}, or \emph{fip}, if the intersection of
every nonempty finite subcollection of \(\mathcal{H}\) is nonempty. That is,
\[
B_1 \cap \cdots \cap B_n \neq \emptyset
\]
for every \(n\) and every \(B_1,\ldots,B_n \in \mathcal{H}\).
\end{defn}

\begin{prop}[Proper filters and the finite intersection property]
\label{prop:proper-filter-iff-fip}
Let \(I\neq\emptyset\), and let \(\mathcal{F}\) be a filter on \(I\).
Then
\[
\mathcal{F}\text{ is proper}
\quad\Longleftrightarrow\quad
\mathcal{F}\text{ has the finite intersection property}.
\]
\end{prop}

\begin{proof}
We first observe that \(\mathcal{F}\) is closed under every nonempty
finite intersection. More precisely, if
\[
A_1,\ldots,A_n\in\mathcal{F},
\qquad n\geq 1,
\]
then
\[
A_1\cap\cdots\cap A_n\in\mathcal{F}.
\]

We prove this by induction on \(n\). For \(n=1\), the assertion is
immediate. Suppose that it holds for some \(n\geq 1\). Then
\[
A_1\cap\cdots\cap A_n\in\mathcal{F}.
\]
Since \(A_{n+1}\in\mathcal{F}\), the intersection axiom gives
\[
(A_1\cap\cdots\cap A_n)\cap A_{n+1}
=
A_1\cap\cdots\cap A_{n+1}
\in\mathcal{F}.
\]
Thus \(\mathcal{F}\) is closed under every nonempty finite intersection.

\noindent\((\Rightarrow)\)
Assume that \(\mathcal{F}\) is proper. Let
\(A_1,\ldots,A_n\in\mathcal{F}\), where \(n\geq 1\). By the observation
above,
\[
A_1\cap\cdots\cap A_n\in\mathcal{F}.
\]
If
\[
A_1\cap\cdots\cap A_n=\emptyset,
\]
then \(\emptyset\in\mathcal{F}\), contradicting the properness of
\(\mathcal{F}\). Therefore,
\[
A_1\cap\cdots\cap A_n\neq\emptyset.
\]
Hence \(\mathcal{F}\) has the finite intersection property.

\noindent\((\Leftarrow)\)
Suppose that \(\mathcal{F}\) has the finite intersection property.
Assume, for contradiction, that \(\mathcal{F}\) is not proper. Then
\[
\emptyset\in\mathcal{F}.
\]
Let \(A\subseteq I\). Since
\[
\emptyset\subseteq A\subseteq I
\qquad\text{and}\qquad
\emptyset\subseteq A^c\subseteq I,
\]
the superset axiom implies that
\[
A\in\mathcal{F}
\qquad\text{and}\qquad
A^c\in\mathcal{F}.
\]
Thus \(\{A,A^c\}\) is a finite subcollection of \(\mathcal{F}\), but
\[
A\cap A^c=\emptyset.
\]
This contradicts the finite intersection property. Therefore,
\(\emptyset\notin\mathcal{F}\), and hence \(\mathcal{F}\) is proper.
\end{proof}

\begin{prop}
The filter \(\mathcal{F}^{\mathcal{H}}\) generated by \(\mathcal{H}\) is
proper if and only if \(\mathcal{H}\) has the finite intersection property.
\end{prop}

\begin{proof}
Recall that
\[
\mathcal{F}^{\mathcal{H}}
=
\left\{
A \subseteq I :
B_1 \cap \cdots \cap B_n \subseteq A
\text{ for some } n \geq 1
\text{ and some } B_1,\ldots,B_n \in \mathcal{H}
\right\}.
\]

\noindent\((\Rightarrow)\)
Assume that \(\mathcal{F}^{\mathcal{H}}\) is proper. Suppose, for a
contradiction, that \(\mathcal{H}\) does not have the finite intersection
property. Then there exist \(B_1,\ldots,B_n \in \mathcal{H}\) such that
\[
B_1 \cap \cdots \cap B_n = \emptyset.
\]
Since
\[
B_1 \cap \cdots \cap B_n \subseteq \emptyset,
\]
the definition of \(\mathcal{F}^{\mathcal{H}}\) gives
\[
\emptyset \in \mathcal{F}^{\mathcal{H}}.
\]
This contradicts the assumption that
\(\mathcal{F}^{\mathcal{H}}\) is proper. Therefore,
\(\mathcal{H}\) has the finite intersection property.

\noindent\((\Leftarrow)\)
Assume that \(\mathcal{H}\) has the finite intersection property. To prove
that \(\mathcal{F}^{\mathcal{H}}\) is proper, it suffices to show that
\[
\emptyset \notin \mathcal{F}^{\mathcal{H}}.
\]

Suppose, for a contradiction, that
\[
\emptyset \in \mathcal{F}^{\mathcal{H}}.
\]
Then, by the definition of \(\mathcal{F}^{\mathcal{H}}\), there exist
\(B_1,\ldots,B_n \in \mathcal{H}\) such that
\[
B_1 \cap \cdots \cap B_n \subseteq \emptyset.
\]
Therefore,
\[
B_1 \cap \cdots \cap B_n = \emptyset.
\]
But this contradicts the finite intersection property of \(\mathcal{H}\).

Hence,
\[
\emptyset \notin \mathcal{F}^{\mathcal{H}},
\]
so \(\mathcal{F}^{\mathcal{H}}\) is proper.
\end{proof}

\begin{prop}[\textit{this used to be an exercise}] \label{prop:adjoining a complement to a proper filter}
Let \(\mathcal{F}\) be a proper filter on \(I\), and let \(A\subseteq I\).
Then
\[
\mathcal{F}\cup\{A^c\}
\]
has the finite intersection property if and only if
\[
A\notin\mathcal{F},
\]
where \(A^c=I\setminus A\).
\end{prop}

\begin{proof}
We prove both implications.

\noindent\((\Rightarrow)\)
Assume that
\[
\mathcal{F}\cup\{A^c\}
\]
has the finite intersection property. Suppose, for the sake of contradiction,
that
\[
A\in\mathcal{F}.
\]
Then both \(A\) and \(A^c\) belong to
\(\mathcal{F}\cup\{A^c\}\). However,
\[
A\cap A^c=\emptyset.
\]
Thus a finite subcollection of
\(\mathcal{F}\cup\{A^c\}\) has empty intersection, contradicting the finite
intersection property. Therefore,
\[
A\notin\mathcal{F}.
\]

\noindent\((\Leftarrow)\)
Assume that
\[
A\notin\mathcal{F}.
\]
We prove that \(\mathcal{F}\cup\{A^c\}\) has the finite intersection
property.

Let \(\mathcal{C}\) be a nonempty finite subcollection of
\(\mathcal{F}\cup\{A^c\}\). There are two cases.

First, suppose that \(A^c\notin\mathcal{C}\). Then every member of
\(\mathcal{C}\) belongs to \(\mathcal{F}\). Since \(\mathcal{F}\) is closed
under finite intersections,
\[
\bigcap_{C\in\mathcal{C}} C\in\mathcal{F}.
\]
Because \(\mathcal{F}\) is proper, it does not contain the empty set.
Therefore,
\[
\bigcap_{C\in\mathcal{C}} C\neq\emptyset.
\]

Now suppose that \(A^c\in\mathcal{C}\). \textcolor{red}{(I have another, lengthy, approach using contradictions, on iPad page 10)} Then the remaining members of
\(\mathcal{C}\), if there are any, belong to \(\mathcal{F}\). Thus we may
write
\[
\bigcap_{C\in\mathcal{C}}C
=
B_1\cap\cdots\cap B_n\cap A^c,
\]
where \(B_1,\ldots,B_n\in\mathcal{F}\). We also allow \(n=0\), in which
case the intersection \(B_1\cap\cdots\cap B_n\) is understood to be \(I\).

Set
\[
B=
\begin{cases}
B_1\cap\cdots\cap B_n, & n\geq 1,\\
I, & n=0.
\end{cases}
\]
Since \(\mathcal{F}\) is a filter, \(I\in\mathcal{F}\), and
\(\mathcal{F}\) is closed under finite intersections. Hence
\[
B\in\mathcal{F}.
\]

Suppose, for the sake of contradiction, that
\[
B\cap A^c=\emptyset.
\]
Then \(B\subseteq A\). Indeed, if \(x\in B\), then \(x\notin A^c\), so
\(x\in A\). Therefore,
\[
B\subseteq A\subseteq I.
\]
Since \(B\in\mathcal{F}\), the superset axiom implies that
\[
A\in\mathcal{F}.
\]
This contradicts our assumption that \(A\notin\mathcal{F}\). Hence
\[
B\cap A^c\neq\emptyset.
\]

Thus every nonempty finite subcollection of
\(\mathcal{F}\cup\{A^c\}\) has nonempty intersection. Therefore,
\[
\mathcal{F}\cup\{A^c\}
\]
has the finite intersection property.
\end{proof}

\subsubsection{Maximal proper filter}

\begin{defn}[Maximal proper filter]
A proper filter \(\mathcal{F}\) on \(I\) is called \emph{maximal} if there
is no proper filter on \(I\) that properly contains \(\mathcal{F}\).

Equivalently, \(\mathcal{F}\) is maximal if, whenever \(\mathcal{G}\) is a
proper filter on \(I\) satisfying
\[
\mathcal{F}\subseteq\mathcal{G},
\]
we necessarily have
\[
\mathcal{F}=\mathcal{G}.
\]
\end{defn}

\begin{prop}[(\textit{this used to be an exercise}) Ultrafilters are maximal proper filters] \label{prop: ultrafilters are maximal proper filters}
Let \(\mathcal{F}\) be a proper filter on \(I\). Then
\(\mathcal{F}\) is an ultrafilter if and only if it is a maximal proper
filter on \(I\).
\end{prop}

\begin{proof}
We prove both implications.

\noindent\((\Rightarrow)\)
Assume that \(\mathcal{F}\) is an ultrafilter. We prove that
\(\mathcal{F}\) is maximal among the proper filters on \(I\).

Let \(\mathcal{G}\) be a proper filter such that
\[
\mathcal{F}\subseteq\mathcal{G}.
\]
We want to show that
\[
\mathcal{G}=\mathcal{F}.
\]

Since \(\mathcal{F}\subseteq\mathcal{G}\), it remains only to prove that
\[
\mathcal{G}\subseteq\mathcal{F}.
\]
Let \(A\in\mathcal{G}\). Because \(\mathcal{F}\) is an ultrafilter, either
\[
A\in\mathcal{F}
\qquad\text{or}\qquad
A^c\in\mathcal{F}.
\]

Suppose that \(A^c\in\mathcal{F}\). Since
\(\mathcal{F}\subseteq\mathcal{G}\), this gives
\[
A^c\in\mathcal{G}.
\]
But \(A\in\mathcal{G}\) as well. Since \(\mathcal{G}\) is closed under
intersections,
\[
A\cap A^c=\emptyset\in\mathcal{G},
\]
contradicting the fact that \(\mathcal{G}\) is proper.

Therefore, the second possibility is impossible, and we must have
\[
A\in\mathcal{F}.
\]
Since \(A\in\mathcal{G}\) was arbitrary,
\[
\mathcal{G}\subseteq\mathcal{F}.
\]
Consequently,
\[
\mathcal{G}=\mathcal{F}.
\]
Thus \(\mathcal{F}\) is a maximal proper filter.

\textcolor{red}{Another way to prove this direction:}
$\call{F}$ ultrafilter. For contradiction: assume there exists $G$ proper filter such that $\call{F} \subseteq \call{G}$. So there exists $A \in \call{G}$ such that $A \notin \call{F}$. Thus $A \notin \call{F}$ implies $A^c \in \call{F}$ which implies $A^c \in \call{G}$ since $\call{F} \subseteq \call{G}$. So $A, A^c \in \call{G}$ implies $\emptyset \in \call{G}$ which is contradiction cos $\call{G}$ is proper.

\noindent\((\Leftarrow)\)
Assume that \(\mathcal{F}\) is a maximal proper filter. We prove that
\(\mathcal{F}\) is an ultrafilter.

Let \(A\subseteq I\). We must show that
\[
A\in\mathcal{F}
\qquad\text{or}\qquad
A^c\in\mathcal{F}.
\]

If \(A\in\mathcal{F}\), there is nothing to prove. Suppose instead that
\[
A\notin\mathcal{F}.
\]
By the preceding proposition,
\[
\mathcal{H}
=
\mathcal{F}\cup\{A^c\}
\]
has the finite intersection property.

Let \(\mathcal{F}^{\mathcal{H}}\) be the filter generated by
\(\mathcal{H}\). Since \(\mathcal{H}\) has the finite intersection
property, the filter \(\mathcal{F}^{\mathcal{H}}\) is proper.

Moreover,
\[
\mathcal{F}\subseteq\mathcal{H}
\subseteq\mathcal{F}^{\mathcal{H}}.
\]
Thus \(\mathcal{F}^{\mathcal{H}}\) is a proper filter containing
\(\mathcal{F}\). Since \(\mathcal{F}\) is maximal among proper filters, we
must have
\[
\mathcal{F}^{\mathcal{H}}=\mathcal{F}.
\]

But
\[
A^c\in\mathcal{H}\subseteq\mathcal{F}^{\mathcal{H}}.
\]
Therefore,
\[
A^c\in\mathcal{F}.
\]

We have proved that for every \(A\subseteq I\),
\[
A\in\mathcal{F}
\qquad\text{or}\qquad
A^c\in\mathcal{F}.
\]
Hence \(\mathcal{F}\) is an ultrafilter.

\textcolor{red}{another way to prove this direction (doubtful):}
Say $\call{F}$ is maximal proper filter. For contradiction, assume there exists $A \subseteq I$ such that $A \notin \call{F}$ and $A^c \notin \call{F}$ (so $\call{F}$ is not ultrafilter). Then $\call{H}=\call{F} \cup \{A^c\}$ has the fip. Then $\call{F}^\call{H}$ is proper. We have $\call{F} \subseteq \call{H} \subseteq \call{F}^\call{H}$. Since $\call{F}$ is maximal, we must have $\call{F}=\call{F}^\call{H}=\call{H}$. This is a contradiction, because $A^c \notin \call{F}$ but $A^c \in \call{H}$. But $\call{F}=\call{H}$. \textcolor{red}{is this proof correct???}
\end{proof}

\begin{prop}
Let \(\mathcal{H} \subseteq \mathcal{P}(I)\) have the finite intersection
property. Then, for every \(A \subseteq I\), at least one of the collections
\[
\mathcal{H} \cup \{A\}
\qquad\text{and}\qquad
\mathcal{H} \cup \{A^c\},
\]
where \(A^c=I\setminus A\), has the finite intersection property.
\end{prop}

\begin{proof}
Let \(A\subseteq I\). We prove the result by contradiction. Suppose that
neither
\[
\mathcal{H}\cup\{A\}
\]
nor
\[
\mathcal{H}\cup\{A^c\}
\]
has the finite intersection property.

Since \(\mathcal{H}\cup\{A\}\) does not have the finite intersection
property, there exist finitely many sets
\[
B_1,\ldots,B_m\in\mathcal{H}
\]
such that
\[
B_1\cap\cdots\cap B_m\cap A=\emptyset.
\]
Likewise, since \(\mathcal{H}\cup\{A^c\}\) does not have the finite
intersection property, there exist finitely many sets
\[
C_1,\ldots,C_n\in\mathcal{H}
\]
such that
\[
C_1\cap\cdots\cap C_n\cap A^c=\emptyset.
\]

Put
\[
B=B_1\cap\cdots\cap B_m
\qquad\text{and}\qquad
C=C_1\cap\cdots\cap C_n.
\]
Then
\[
B\cap A=\emptyset
\qquad\text{and}\qquad
C\cap A^c=\emptyset.
\]

The first equality implies that \(B\subseteq A^c\). Indeed, if \(x\in B\),
then \(x\notin A\), since otherwise \(x\in B\cap A\), contradicting
\(B\cap A=\emptyset\). Hence \(x\in A^c\).

Similarly, the second equality implies that \(C\subseteq A\). Therefore,
\[
B\cap C
\subseteq
A^c\cap A
=
\emptyset.
\]
Thus,
\[
B\cap C=\emptyset.
\]

But
\[
B\cap C
=
B_1\cap\cdots\cap B_m\cap C_1\cap\cdots\cap C_n,
\]
which is a finite intersection of members of \(\mathcal{H}\). This
contradicts the assumption that \(\mathcal{H}\) has the finite intersection
property.

Therefore, at least one of
\[
\mathcal{H}\cup\{A\}
\qquad\text{or}\qquad
\mathcal{H}\cup\{A^c\}
\]
has the finite intersection property.
\end{proof}

\subsection{Zorn's Lemma}
\label{subsec:zorns-lemma}

The preceding proposition suggests a possible way of constructing an
ultrafilter. We could begin with a collection
\[
\mathcal{H}\subseteq\mathcal{P}(I)
\]
having the finite intersection property, and then consider every subset
\(A\subseteq I\), one at a time. At each step, we would add either \(A\) or
its complement \(A^c\), choosing one that preserves the finite intersection
property.

If this process could be continued until every subset of \(I\) had been
considered, the resulting collection would decide, for every
\(A\subseteq I\), whether \(A\) or \(A^c\) should count as large. We would
therefore expect the resulting object to be an ultrafilter.

The difficulty lies in the phrase ``consider every subset one at a time.''
The power set \(\mathcal{P}(I)\) may be extremely large, and there may be no
ordinary sequence
\[
A_1,A_2,A_3,\ldots
\]
containing all of its members. To carry out this procedure directly, one
would need to arrange \(\mathcal{P}(I)\) in a suitable well-ordered list and
continue the construction through possibly transfinite stages.

The statement that every set can be well-ordered is equivalent to the
\emph{axiom of choice}. The form of the axiom of choice most commonly used
in algebra is \emph{Zorn's lemma}. It allows us to prove that a maximal
object exists without explicitly constructing a transfinite list of all the
possible extensions.

\subsubsection{Partially ordered sets}

\begin{defn}[Partially ordered set]
A \emph{partially ordered set}, or \emph{poset}, is a pair
\[
(P,\leq),
\]
where \(P\) is a set and \(\leq\) is a relation on \(P\) satisfying:
\begin{enumerate}
    \item \(p\leq p\) for every \(p\in P\);
    \item if \(p\leq q\) and \(q\leq p\), then \(p=q\);
    \item if \(p\leq q\) and \(q\leq r\), then \(p\leq r\).
\end{enumerate}

The elements of a partially ordered set need not all be comparable. It may
happen that neither \(p\leq q\) nor \(q\leq p\).
\end{defn}

In the application to filters, the elements of \(P\) will themselves be
filters, and the order relation will be set inclusion:
\[
\mathcal{F}\leq\mathcal{G}
\quad\Longleftrightarrow\quad
\mathcal{F}\subseteq\mathcal{G}.
\]
Thus, moving upward in the partial order means enlarging a filter by
declaring more subsets of \(I\) to be large.

\begin{defn}[Chain]
Let \((P,\leq)\) be a partially ordered set. A subset
\(\mathcal{C}\subseteq P\) is called a \emph{chain} if every two of its
members are comparable. In other words, for every \(p,q\in\mathcal{C}\),
either
\[
p\leq q
\qquad\text{or}\qquad
q\leq p.
\]
\end{defn}

For filters ordered by inclusion, a chain is therefore a collection of
filters such that any two are related by inclusion. A typical chain may be
pictured as
\[
\mathcal{F}_1
\subseteq
\mathcal{F}_2
\subseteq
\mathcal{F}_3
\subseteq
\cdots.
\]
The chain represents a collection of mutually compatible stages in a
process of enlargement.

\begin{defn}[Upper bound]
Let \(\mathcal{C}\) be a chain in a partially ordered set \(P\). An element
\(u\in P\) is an \emph{upper bound} for \(\mathcal{C}\) if
\[
p\leq u
\]
for every \(p\in\mathcal{C}\).
\end{defn}

An upper bound need not belong to the chain itself. It is simply an element
of \(P\) that lies above every member of the chain.

For a chain of filters ordered by inclusion, the natural candidate for an
upper bound is their union:
\[
\bigcup_{\mathcal{F}\in\mathcal{C}}\mathcal{F}.
\]
This union combines all the sets that were declared large at any stage of
the chain.

\begin{defn}[Maximal element]
Let \((P,\leq)\) be a partially ordered set. An element \(m\in P\) is
\emph{maximal} if there is no \(p\in P\) such that
\[
m\leq p
\qquad\text{and}\qquad
m\neq p.
\]
Equivalently, \(m\) is maximal if it cannot be strictly enlarged within
\(P\).
\end{defn}

A maximal element should not be confused with a \emph{maximum} element. A
maximum element is greater than every element of \(P\), whereas a maximal
element is merely an element above which no strictly larger element exists.
A partially ordered set may have several maximal elements and no maximum
element.

In our application, a maximal element will be a proper filter that cannot
be enlarged while remaining proper.

\begin{prop}[Zorn's lemma]
\label{prop:zorns-lemma}
Let \((P,\leq)\) be a partially ordered set. Suppose that every chain in
\(P\) has an upper bound in \(P\). Then \(P\) contains a maximal element.
\end{prop}

The intuition behind Zorn's lemma is the following:
\[
\boxed{
\begin{minipage}{0.82\textwidth}
If every compatible process of enlargement has a valid limiting stage,
then there is an object that cannot be enlarged any further.
\end{minipage}
}
\]

Indeed, imagine beginning with some \(p_0\in P\). If \(p_0\) is not
maximal, choose a larger element \(p_1\). If \(p_1\) is not maximal, choose
\(p_2\), and continue:
\[
p_0<p_1<p_2<\cdots.
\]
These elements form a chain. By the hypothesis of Zorn's lemma, the chain
has an upper bound, which we may denote informally by \(p_\omega\):
\[
p_0<p_1<\cdots<p_n<\cdots\leq p_\omega.
\]
If \(p_\omega\) is not maximal, we continue above it. The process may have
to continue through stages beyond the natural numbers.

Informally, this construction cannot continue forever without eventually
running out of new elements of \(P\). Making this statement precise
requires ordinal numbers, well-orderings, and transfinite recursion.
Zorn's lemma allows us to use the conclusion of that machinery directly,
without rebuilding it in each application.

\subsubsection{Extending a filter to an ultrafilter}

We now apply Zorn's lemma to filters. The central idea is to consider all
proper filters satisfying our original requirements and order them by
inclusion. Zorn's lemma will then produce a filter that cannot be enlarged
any further.

\begin{thm}[Ultrafilter extension theorem]
\label{thm:ultrafilter-extension}
Let \(I\) be a nonempty set, and let
\[
\mathcal{H}\subseteq\mathcal{P}(I)
\]
have the finite intersection property. Then there exists an ultrafilter
\(\mathcal{U}\) on \(I\) such that
\[
\mathcal{H}\subseteq\mathcal{U}.
\]
In other words, every collection having the finite intersection property
can be extended to an ultrafilter.
\end{thm}

\begin{proof}
Let $\mathcal{F}^{\mathcal{H}}$ be the filter generated by \(\mathcal{H}\). Since \(\mathcal{H}\) has the finite intersection property, the previously established characterization of generated filters tells us that
$\mathcal{F}^{\mathcal{H}}$
is proper.

Now let \(P\) be the collection of all proper filters on \(I\) that contain
\(\mathcal{F}^{\mathcal{H}}\):
\[
P
=
\left\{
\mathcal{F}:
\mathcal{F}\text{ is a proper filter on }I
\text{ and }
\mathcal{F}^{\mathcal{H}}\subseteq\mathcal{F}
\right\}.
\]
We order \(P\) by set inclusion.

The collection \(P\) is nonempty because
\[
\mathcal{F}^{\mathcal{H}}\in P.
\]

We now verify the hypothesis of Zorn's lemma. Let
\[
\mathcal{C}\subseteq P
\]
be a chain. If \(\mathcal{C}=\emptyset\), then any member of \(P\) is an
upper bound, so suppose that \(\mathcal{C}\neq\emptyset\).

Define
\[
\mathcal{G}
=
\bigcup_{\mathcal{F}\in\mathcal{C}}\mathcal{F}.
\]

Because \(\mathcal{C}\) is linearly ordered by inclusion,
Proposition~\ref{prop:union of a linearly ordered collection of filters is a filter}
shows that \(\mathcal{G}\) is a filter on \(I\).

Moreover, every \(\mathcal{F}\in\mathcal{C}\) contains
\(\mathcal{F}^{\mathcal{H}}\). Therefore,
\[
\mathcal{F}^{\mathcal{H}}
\subseteq
\mathcal{G}.
\]

The filter \(\mathcal{G}\) is also proper. Indeed, if
\[
\emptyset\in\mathcal{G},
\]
then by the definition of the union there would be some
\(\mathcal{F}\in\mathcal{C}\) such that
\[
\emptyset\in\mathcal{F}.
\]
This is impossible because every member of \(\mathcal{C}\) belongs to \(P\)
and is therefore a proper filter.

Thus,
\[
\mathcal{G}\in P.
\]
Since every \(\mathcal{F}\in\mathcal{C}\) satisfies
\[
\mathcal{F}\subseteq\mathcal{G},
\]
the filter \(\mathcal{G}\) is an upper bound for the chain
\(\mathcal{C}\).

We have shown that every chain in \(P\) has an upper bound in \(P\).
Therefore, by Zorn's lemma, \(P\) has a maximal element. Denote it by
\(\mathcal{U}\).

Thus,
\[
\mathcal{F}^{\mathcal{H}}
\subseteq
\mathcal{U},
\]
and \(\mathcal{U}\) is a proper filter that cannot be properly enlarged
while remaining in \(P\).

In fact, \(\mathcal{U}\) is maximal among all proper filters on \(I\).
Suppose that \(\mathcal{G}\) is a proper filter satisfying
\[
\mathcal{U}\subseteq\mathcal{G}.
\]
Since
\[
\mathcal{F}^{\mathcal{H}}
\subseteq
\mathcal{U}
\subseteq
\mathcal{G},
\]
we have \(\mathcal{G}\in P\). The maximality of \(\mathcal{U}\) in \(P\)
therefore gives
\[
\mathcal{G}=\mathcal{U}.
\]

Consequently, \(\mathcal{U}\) is a maximal proper filter. By
Proposition~\ref{prop: ultrafilters are maximal proper filters},
\(\mathcal{U}\) is an ultrafilter.

Finally,
\[
\mathcal{H}
\subseteq
\mathcal{F}^{\mathcal{H}}
\subseteq
\mathcal{U},
\]
so \(\mathcal{U}\) is an ultrafilter extending the original collection
\(\mathcal{H}\).
\end{proof}

The logical structure of the proof can be summarized as follows:
\[
\boxed{
\begin{gathered}
\mathcal{H}\text{ has the finite intersection property}
\\[2mm]
\Downarrow
\\[2mm]
\mathcal{H}\text{ generates a proper filter }
\mathcal{F}^{\mathcal{H}}
\\[2mm]
\Downarrow
\\[2mm]
\text{consider all proper filters extending }
\mathcal{F}^{\mathcal{H}}
\\[2mm]
\Downarrow
\\[2mm]
\text{unions of chains provide upper bounds}
\\[2mm]
\Downarrow\quad\text{Zorn's lemma}
\\[2mm]
\text{there exists a maximal proper filter}
\\[2mm]
\Downarrow
\\[2mm]
\text{there exists an ultrafilter extending }\mathcal{H}.
\end{gathered}
}
\]

Thus the theorem is a completion principle:
\[
\boxed{
\text{finite consistency}
\quad\Longrightarrow\quad
\text{a maximal consistent notion of largeness}.
}
\]
The finite intersection property guarantees that the initial requirements
do not contradict one another. Zorn's lemma guarantees that these
requirements can be enlarged until every subset of \(I\) has been decided.

A result that follows immediately is the following:

\prop[\textit{This used to be an exercise}]{Let $A \subseteq I$ be nonempty, $A \neq \emptyset$. Then there is an ultrafilter $\call{U}$ on $I$ with $A \in \call{U}$.}
\begin{proof}
    Let \(\emptyset\neq A\subseteq I\), and set
\[
\mathcal{H}=\{A\}.
\]
Since \(A\neq\emptyset\), the filter \(\mathcal{F}^{\mathcal{H}}\) generated by
\(\mathcal{H}\) is proper. Hence \(\mathcal{H}\) has the finite intersection
property.

By the ultrafilter extension theorem, there exists an ultrafilter
\(\mathcal{U}\) on \(I\) such that
\[
\mathcal{H}\subseteq\mathcal{U}.
\]
Since \(A\in\mathcal{H}\), it follows that
\[
A\in\mathcal{U}.
\]
\end{proof}

\subsubsection{Existence of nonprincipal ultrafilters}

The ultrafilter extension theorem becomes especially important when the
underlying set \(I\) is infinite.

\begin{cor}[Existence of nonprincipal ultrafilters]
\label{cor:nonprincipal-ultrafilter}
Every infinite set admits a nonprincipal ultrafilter.
\end{cor}

\begin{proof}
Let \(I\) be infinite, and consider the cofinite filter
\[
\mathcal{F}^{\mathrm{co}}
=
\left\{
A\subseteq I:
I\setminus A\text{ is finite}
\right\}.
\]
Since \(I\) is infinite, the cofinite filter is proper.

Every proper filter has the finite intersection property: if
\[
A_1,\ldots,A_n\in\mathcal{F}^{\mathrm{co}},
\]
then
\[
A_1\cap\cdots\cap A_n\in\mathcal{F}^{\mathrm{co}}.
\]
Because the filter is proper, this intersection cannot be empty.

We may therefore apply
Theorem~\ref{thm:ultrafilter-extension} to obtain an ultrafilter
\(\mathcal{U}\) such that
\[
\mathcal{F}^{\mathrm{co}}
\subseteq
\mathcal{U}.
\]

We show that \(\mathcal{U}\) is not principal. Let \(i\in I\). The set
\[
I\setminus\{i\}
\]
is cofinite, so
\[
I\setminus\{i\}
\in
\mathcal{F}^{\mathrm{co}}
\subseteq
\mathcal{U}.
\]
It follows that
\[
\{i\}\notin\mathcal{U}.
\]
Indeed, if both \(\{i\}\) and \(I\setminus\{i\}\) belonged to
\(\mathcal{U}\), then their intersection would also belong to
\(\mathcal{U}\), giving
\[
\emptyset
=
\{i\}\cap\bigl(I\setminus\{i\}\bigr)
\in\mathcal{U},
\]
contrary to the properness of \(\mathcal{U}\).

Thus \(\mathcal{U}\) contains no singleton. But the principal ultrafilter
generated by \(i\), described in
Example~\ref{eg:principal ultrafilter generated by i}, contains the
singleton \(\{i\}\). Therefore, \(\mathcal{U}\) cannot be principal.
\end{proof}

A principal ultrafilter declares a set large when it contains one particular
point. A nonprincipal ultrafilter does not concentrate on any individual
point. It contains every cofinite subset of \(I\), contains no finite
subset, and nevertheless decides every subset \(A\subseteq I\):
\[
A\in\mathcal{U}
\qquad\text{or}\qquad
A^c\in\mathcal{U}.
\]
It therefore gives an abstract and highly decisive notion of largeness
without selecting a distinguished element of \(I\).

\paragraph{Connection with the construction of the hyperreal numbers.}
The existence of a nonprincipal ultrafilter is the key result needed for
the ultrapower construction of the hyperreal number system. Taking
\(I=\mathbb{N}\), Corollary~\ref{cor:nonprincipal-ultrafilter} allows us to
fix a nonprincipal ultrafilter \(\mathcal{U}\) on \(\mathbb{N}\). For two
real-valued sequences
\[
r=(r_n)_{n\in\mathbb{N}}
\qquad\text{and}\qquad
s=(s_n)_{n\in\mathbb{N}},
\]
we will declare
\[
r\sim_{\mathcal{U}}s
\quad\Longleftrightarrow\quad
\{n\in\mathbb{N}:r_n=s_n\}\in\mathcal{U}.
\]
The hyperreal numbers will then be obtained from equivalence classes of
real sequences:
\[
{}^{\ast}\mathbb{R}
=
\mathbb{R}^{\mathbb{N}}/{\sim_{\mathcal{U}}}.
\]

The fact that \(\mathcal{U}\) is an ultrafilter ensures that it makes a
consistent decision about every subset of \(\mathbb{N}\). This decisiveness
will be essential when defining and comparing hyperreal numbers. The fact
that \(\mathcal{U}\) is nonprincipal ensures that every cofinite set counts
as large while no single index is distinguished.

For example, the sequence
\[
\left(\frac{1}{n}\right)_{n\geq 1}
\]
is not equivalent to the zero sequence, because
\[
\left\{
n\in\mathbb{N}:
\frac{1}{n}=0
\right\}
=
\emptyset
\notin\mathcal{U}.
\]
However, for every \(m\geq 1\), the set
\[
\left\{
n\in\mathbb{N}:
0<\frac{1}{n}<\frac{1}{m}
\right\}
\]
contains all \(n>m\) and is therefore cofinite. Hence it belongs to
\(\mathcal{U}\). The resulting hyperreal is consequently nonzero but
smaller than every positive real number of the form \(1/m\): it is an
infinitesimal.

A principal ultrafilter would merely focus on one fixed term of each
sequence and would not produce genuinely new numbers. The nonprincipal
ultrafilter supplied by Zorn's lemma is therefore what makes the
construction of nonzero infinitesimals, infinitely large numbers, and the
full hyperreal extension possible.



\subsection{Summary of Definitions and Main Results on Filters}
\label{subsec:filters-summary}

Throughout, let \(I\neq\emptyset\), let
\(\mathcal{P}(I)=\{A:A\subseteq I\}\), and write
\(A^c=I\setminus A\). The following statements collect the definitions and
results established above, in the order in which they are used.

\subsubsection*{1. Basic filter structure}

\begin{description}
    \item[\textnormal{(F1)}] \textbf{Filter.}
    A filter on \(I\) is a nonempty collection
    \(\mathcal{F}\subseteq\mathcal{P}(I)\) such that
    \[\text{(intersection axiom)} \qquad
    A,B\in\mathcal{F}
    \Longrightarrow
    A\cap B\in\mathcal{F},
    \]
    and
    \[ \text{(superset axiom)} \qquad
    A\in\mathcal{F},\quad A\subseteq B\subseteq I
    \Longrightarrow
    B\in\mathcal{F}.
    \]
    Consequently, \(I\in\mathcal{F}\), and \(\mathcal{F}\) is closed under
    every nonempty finite intersection.

    \item[\textnormal{(F2)}] \textbf{Nonemptiness criterion.}
    If \(\mathcal{F}\subseteq\mathcal{P}(I)\) satisfies the superset axiom,
    then
    \[
    \mathcal{F}\neq\emptyset
    \quad\Longleftrightarrow\quad
    I\in\mathcal{F}.
    \]

    \item[\textnormal{(F3)}] \textbf{The improper filter.}
    For every filter \(\mathcal{F}\) on \(I\),
    \[
    \emptyset\in\mathcal{F}
    \quad\Longleftrightarrow\quad
    \mathcal{F}=\mathcal{P}(I).
    \]
    A filter is \emph{proper} precisely when
    \(\emptyset\notin\mathcal{F}\).
\end{description}

\subsubsection*{2. Ultrafilters and their Boolean calculus}

\begin{description}
    \item[\textnormal{(U1)}] \textbf{Ultrafilter.}
    An ultrafilter \(\mathcal{U}\) on \(I\) is a proper filter such that, for
    every \(A\subseteq I\),
    \[
    A\in\mathcal{U}
    \qquad\text{or}\qquad
    A^c\in\mathcal{U}.
    \]
    Since \(\mathcal{U}\) is proper, these alternatives are mutually
    exclusive. Thus exactly one of \(A\) and \(A^c\) belongs to
    \(\mathcal{U}\).

    \item[\textnormal{(U2)}] \textbf{Intersection, union, and complement laws.}
    For \(A,B\subseteq I\), an ultrafilter $\call{U}$ satisfies
    \[
    A\cap B\in\mathcal{U}
    \quad\Longleftrightarrow\quad
    A\in\mathcal{U}\ \text{and}\ B\in\mathcal{U},
    \]
    \[
    A\cup B\in\mathcal{U}
    \quad\Longleftrightarrow\quad
    A\in\mathcal{U}\ \text{or}\ B\in\mathcal{U},
    \]
    and
    \[
    A^c\in\mathcal{U}
    \quad\Longleftrightarrow\quad
    A\notin\mathcal{U}.
    \]
    More generally, for \(A_1,\ldots,A_n\subseteq I\),
    \[
    \bigcap_{k=1}^{n}A_k\in\mathcal{U}
    \quad\Longleftrightarrow\quad
    A_k\in\mathcal{U}\ \text{for every }k,
    \]
    and
    \[
    \bigcup_{k=1}^{n}A_k\in\mathcal{U}
    \quad\Longleftrightarrow\quad
    A_k\in\mathcal{U}\ \text{for at least one }k.
    \]

    \item[\textnormal{(U3)}] \textbf{Finite disjoint unions.}
    If \(A_1,\ldots,A_n\subseteq I\) are pairwise disjoint ($A_i \cap A_j = \emptyset, i\neq j$) and
    \[
    A_1\cup\cdots\cup A_n\in\mathcal{U},
    \]
    then exactly one of the sets \(A_1,\ldots,A_n\) belongs to
    \(\mathcal{U}\).
\end{description}

\subsubsection*{3. Canonical examples and principal ultrafilters}

\begin{description}
    \item[\textnormal{(E1)}] \textbf{Principal ultrafilter generated by a point.}
    For \(i\in I\),
    \[
    \mathcal{F}^{i}
    =
    \{A\subseteq I:i\in A\}
    =
    \{A\subseteq I:\{i\}\subseteq A\}
    \]
    is a proper ultrafilter. A set is large with respect to
    \(\mathcal{F}^{i}\) exactly when it contains \(i\).

    \item[\textnormal{(E2)}]
\textbf{Singleton criterion for principal ultrafilters
(derived in Proposition~\ref{prop:finite-member-implies-principal}).}
If \(\mathcal{U}\) is an ultrafilter and \(\{i\}\in\mathcal{U}\), then
\[
\mathcal{U}=\mathcal{F}^{i}.
\]
Indeed, for every \(A\subseteq I\),
\[
A\in\mathcal{U}
\quad\Longleftrightarrow\quad
i\in A.
\]
Consequently, if an ultrafilter contains a finite set, then it contains a
singleton and is principal. In particular, every ultrafilter on a finite
    set is principal.

    \item[\textnormal{(E3)}] \textbf{Nonprincipal ultrafilters.}
    A nonprincipal ultrafilter contains no finite set and contains every
    cofinite set.

    \item[\textnormal{(E4)}] \textbf{Cofinite filter.}
    The collection
    \[
    \mathcal{F}^{\mathrm{co}}
    =
    \{A\subseteq I:I\setminus A\text{ is finite}\}
    \]
    is a filter. It is proper if and only if \(I\) is infinite, and it is
    not an ultrafilter. When \(I\) is infinite, there exist sets
    \(A\subseteq I\) for which both \(A\) and \(A^c\) are infinite, so
    neither belongs to \(\mathcal{F}^{\mathrm{co}}\).
\end{description}

\subsubsection*{4. Generated filters and finite consistency}

\begin{description}
    \item[\textnormal{(G1)}] \textbf{Filter generated by a collection.}
    Let \(\emptyset\neq\mathcal{H}\subseteq\mathcal{P}(I)\). The filter
    generated by \(\mathcal{H}\) is
    \[
    \mathcal{F}^{\mathcal{H}}
    =
    \left\{
    A\subseteq I:
    B_1\cap\cdots\cap B_n\subseteq A
    \text{ for some }n\geq 1
    \text{ and }B_1,\ldots,B_n\in\mathcal{H}
    \right\}.
    \]
    It contains \(\mathcal{H}\), is a filter, and is the smallest filter
    containing \(\mathcal{H}\): for every filter \(\mathcal{G}\),
    \[
    \mathcal{H}\subseteq\mathcal{G}
    \quad\Longrightarrow\quad
    \mathcal{F}^{\mathcal{H}}\subseteq\mathcal{G}.
    \]

    \item[\textnormal{(G2)}] \textbf{Principal filter generated by a set.}
    If \(\mathcal{H}=\{B\}\), then
    \[
    \mathcal{F}^{\mathcal{H}}
    =
    \{A\subseteq I:B\subseteq A\}.
    \]
    This is the principal filter generated by \(B\). The principal
    ultrafilter \(\mathcal{F}^{i}\) is the special case \(B=\{i\}\).

    \item[\textnormal{(G3)}] \textbf{Properness of a principal filter.}
If \(\mathcal{H}=\{B\}\), then
\[
\mathcal{F}^{\mathcal{H}} \text{ is proper}
\quad\Longleftrightarrow\quad
B\neq\varnothing.
\]

    \item[\textnormal{(G4)}] \textbf{Finite intersection property.}
    A collection \(\mathcal{H}\subseteq\mathcal{P}(I)\) has the
    \emph{finite intersection property} (fip) if
    \[
    B_1\cap\cdots\cap B_n\neq\emptyset
    \]
    for every \(n\geq 1\) and every
    \(B_1,\ldots,B_n\in\mathcal{H}\).

    \item[\textnormal{(G5)}]
\textbf{Proper filters and the finite intersection property.}
If \(\mathcal{F}\) is a filter on \(I\), then
\[
\mathcal{F}\text{ is proper}
\quad\Longleftrightarrow\quad
\mathcal{F}\text{ has the finite intersection property}.
\]

    \item[\textnormal{(G6)}] \textbf{Properness of a generated filter.}
    \[
    \mathcal{F}^{\mathcal{H}}\text{ is proper}
    \quad\Longleftrightarrow\quad
    \mathcal{H}\text{ has the finite intersection property}.
    \]

    \item[\textnormal{(G7)}] \textbf{Adjoining a complement to a proper filter.}
    If \(\mathcal{F}\) is a proper filter and \(A\subseteq I\), then
    \[
    \mathcal{F}\cup\{A^c\}\text{ has the fip}
    \quad\Longleftrightarrow\quad
    A\notin\mathcal{F}.
    \]
    Equivalently,
    \[
    \mathcal{F}\cup\{A\}\text{ has the fip}
    \quad\Longleftrightarrow\quad
    A^c\notin\mathcal{F}.
    \]

    \item[\textnormal{(G8)}] \textbf{One side can always be adjoined.}
    If \(\mathcal{H}\) has the fip, then for every \(A\subseteq I\), at
    least one of
    \[
    \mathcal{H}\cup\{A\}
    \qquad\text{and}\qquad
    \mathcal{H}\cup\{A^c\}
    \]
    has the fip.
\end{description}

\subsubsection*{5. Maximal filters and ordered families}

\begin{description}
    \item[\textnormal{(M1)}] \textbf{Maximal proper filter.}
    A proper filter \(\mathcal{F}\) is maximal if no proper filter properly
    contains it; equivalently, whenever \(\mathcal{G}\) is proper and
    \(\mathcal{F}\subseteq\mathcal{G}\), one has
    \(\mathcal{F}=\mathcal{G}\).

    \item[\textnormal{(M2)}] \textbf{Maximality characterization.}
    For every proper filter \(\mathcal{F}\),
    \[
    \mathcal{F}\text{ is an ultrafilter}
    \quad\Longleftrightarrow\quad
    \mathcal{F}\text{ is a maximal proper filter}.
    \]

    \item[\textnormal{(M3)}] \textbf{Partially ordered set.}
    A poset is a pair \((P,\leq)\) in which \(\leq\) is reflexive,
    antisymmetric, and transitive. For filters, the relevant order is
    inclusion:
    \[
    \mathcal{F}\leq\mathcal{G}
    \quad\Longleftrightarrow\quad
    \mathcal{F}\subseteq\mathcal{G}.
    \]

    \item[\textnormal{(M4)}] \textbf{Chains, upper bounds, and maximal elements.}
    A chain \(\mathcal{C}\subseteq P\) is a subset whose members are
    pairwise comparable. An element \(u\in P\) is an upper bound for
    \(\mathcal{C}\) if \(p\leq u\) for every \(p\in\mathcal{C}\). An
    element \(m\in P\) is maximal if no strictly larger element of \(P\)
    lies above it. A maximal element need not be a maximum element.

    \item[\textnormal{(M5)}] \textbf{Union of a chain of filters.}
    If \(\emptyset\neq\mathcal{C}\) is a collection of filters on \(I\)
    linearly ordered by inclusion, then
    \[
    \bigcup_{\mathcal{F}\in\mathcal{C}}\mathcal{F}
    \]
    is a filter on \(I\) and is an upper bound for \(\mathcal{C}\). If
    every member of \(\mathcal{C}\) is proper, then this union is proper.
\end{description}

\subsubsection*{6. Zorn's lemma and existence results}

\begin{description}
    \item[\textnormal{(Z1)}] \textbf{Zorn's lemma.}
    Let \((P,\leq)\) be a partially ordered set. If every chain in \(P\)
    has an upper bound in \(P\), then \(P\) contains a maximal element.

    \item[\textnormal{(Z2)}] \textbf{Ultrafilter extension theorem.}
    If \(\mathcal{H}\subseteq\mathcal{P}(I)\) has the finite intersection
    property, then there exists an ultrafilter \(\mathcal{U}\) on \(I\)
    such that
    \[
    \mathcal{H}\subseteq\mathcal{U}.
    \]

    \item[\textnormal{(Z3)}] \textbf{For any nonempty set there is an ultrafilter that contains that set.}
    If $\emptyset \neq A \subseteq I$, there is ultrafilter $\call{U}$ on $I$ such that $A \in \call{U}$.

    \item[\textnormal{(Z4)}]
\textbf{Corollary of the ultrafilter extension theorem.}
Every proper filter \(\mathcal{F}\) on \(I\) is contained in an ultrafilter.
Indeed, if \(A_1,\ldots,A_n\in\mathcal{F}\), then
\[
A_1\cap\cdots\cap A_n\in\mathcal{F}.
\]
Since \(\mathcal{F}\) is proper, this intersection cannot be empty. Hence
\(\mathcal{F}\) has the finite intersection property. Applying the
ultrafilter extension theorem with \(\mathcal{H}=\mathcal{F}\) yields an
ultrafilter \(\mathcal{U}\) such that
\[
\mathcal{F}\subseteq\mathcal{U}.
\]


    \item[\textnormal{(Z5)}] \textbf{Existence of nonprincipal ultrafilters.}
    Every infinite set \(I\) admits a nonprincipal ultrafilter. More
    precisely, the proper cofinite filter \(\mathcal{F}^{\mathrm{co}}\) can
    be extended to an ultrafilter \(\mathcal{U}\), and any such extension
    is nonprincipal:
    \[
    \mathcal{F}^{\mathrm{co}}
    \subseteq
    \mathcal{U}.
    \]
\end{description}

\subsubsection*{Logical dependency of the main existence theorem}

\[
\begin{array}{c}
\mathcal{H}\text{ has the finite intersection property}
\\[1mm]
\Downarrow
\\[1mm]
\mathcal{F}^{\mathcal{H}}\text{ is a proper filter}
\\[1mm]
\Downarrow
\\[1mm]
\text{consider all proper filters extending }\mathcal{F}^{\mathcal{H}}
\\[1mm]
\Downarrow
\\[1mm]
\text{unions of chains give proper upper bounds}
\\[1mm]
\Downarrow\ \text{Zorn's lemma}
\\[1mm]
\text{a maximal proper filter exists}
\\[1mm]
\Downarrow
\\[1mm]
\text{an ultrafilter extending }\mathcal{H}\text{ exists.}
\end{array}
\]

\subsection{Exercises on Filters}
\prop{If $\emptyset \neq A \subseteq I$, there is an ultrafilter $\call{U}$ on $I$ with $A \in \call{U}$}. 
\begin{proof}
Suppose that \(\emptyset \neq A \subseteq I\). Let
\[
\mathcal{H}=\{A\}.
\]
Then \(\mathcal{H}\) has the finite intersection property, since
\[
\bigcap \mathcal{H}=A\neq\emptyset.
\]
Hence, by the Ultrafilter Extension Theorem (Theorem \ref{thm:ultrafilter-extension}), there exists an ultrafilter
\(\mathcal{U}\) on \(I\) such that
\[
\mathcal{H}\subseteq\mathcal{U}.
\]
Since \(A\in\mathcal{H}\), it follows that
\[
A\in\mathcal{U}.
\]
\end{proof}

\prop{There exists a nonprincipal ultrafilter on $\Ns$ containing the set of even numbers, and another containing the set of odd numbers.}
\begin{proof}
    Let us consider the cofinite filter $\call{F}^{\mathrm{co}}$ on $\Ns$. Since $\Ns$ is infinite, $\call{F}^{\mathrm{co}}$ is proper. Define 
    $$
    E=\{2n: n\in \Ns\} \qquad \text{and} \qquad O =\{2n+1 :n\in \Ns\} 
    $$
    to be the set of the even and odd numbers, respectively. It follows that they are the complement of each other, i.e. $E=\Ns \setminus O$ and $O = \Ns \setminus E$, both are infinite, thus neither of them belongs to $\call{F}^{\mathrm{co}}$. By proposition \ref{prop:adjoining a complement to a proper filter}, since 
    $$
    E \notin \call{F}^{\mathrm{co}} \qquad \text{and} \qquad O \notin \call{F}^{\mathrm{co}}
    $$
    we must have that both 
    $$
    \call{H} = \call{F}^{\mathrm{co}} \cup \{E\} \qquad \text{and} \qquad \call{G}=\call{F}^{\mathrm{co}} \cup \{O\}
    $$
    have the finite intersection property. Thus Theorem \ref{thm:ultrafilter-extension} implies that there exist ultrafilters $\call{U}$ and $\call{U'}$ such that
    $$
    \call{H} \subseteq \call{U} \qquad \text{and} \qquad \call{G} \subseteq \call{U'}
    $$
    where $\call{U}$ is an ultrafilter containing the even numbers and $\call{U'}$ is an ultrafilter containing the odd numbers.
\end{proof}

\prop{An ultrafilter on a finite set must be principal}
\begin{proof}
    Let \(\mathcal{U}\) be an ultrafilter on the finite set
\[
I = \{x_1,\ldots,x_n\}.
\]
Since \(I \in \mathcal{U}\) and \(I\) is finite, Proposition \ref{prop:finite-member-implies-principal} implies that
\(\mathcal{U}\) contains a singleton set \(\{x_i\}\) for some
\(i \in \{1,\ldots,n\}\). Hence \(\mathcal{U}\) is principal, namely
\[
\mathcal{U} = \{A \subseteq I : x_i \in A\}.
\]
\end{proof}
